% Part: sets-functions-relations
% Chapter: infinite
% Section: dedekinds-proof
%
\documentclass[../../../include/open-logic-section]{subfiles}

\begin{document}

\olfileid{sfr}{infinite}{dedekindsproof}

\olsection[डेडेकिंड यांची ``सिद्धता'']{अनंत संचाच्या अस्तित्वाची
डेडेकिंड यांची ``सिद्धता''}

या प्रकरणात आपण डेडेकिंड बीजसंरचनांच्या साहाय्याने नैसर्गिक संख्यांचे
संचसैद्धान्तिक विवेचन दिले. \olref[arith][ref]{sec} मध्ये आपण इतर
गोष्टींबरोबर विश्लेषणाच्या अंकगणितीकरणाच्या तात्त्विक महत्त्वाचा विचार
केला. आता येथे आपण जे साध्य केले आहे त्याच्या महत्त्वाचा विचार करायला हवा.

संपूर्ण \olref[sfr][arith][]{chap} मध्ये आपण नैसर्गिक संख्या आयत्या
मानल्या आणि त्यांचा वापर करून पूर्णांक, परिमेय संख्या आणि वास्तव संख्या
यांची स्पष्ट रचना केली. या प्रकरणात आपण नैसर्गिक संख्यांची स्पष्ट रचना
दिलेली नाही. आपण एवढेच दाखवले आहे की \emph{कोणताही डेडेकिंड-अनंत संच
दिला असता}, आपण असा संच परिभाषित करू शकतो जो आपल्याला~$\Nat$ ने जसे
वागावेसे वाटते अगदी तसाच वागेल.

त्यामुळे \emph{कोणत्या वस्तू नैसर्गिक संख्या आहेत} अशा सत्तामीमांसक
प्रश्नाचे उत्तर आपण दिले आहे, असा दावा अर्थातच करता येणार नाही. पण ही
चांगलीच गोष्ट आहे. शेवटी, \olref[sfr][arith][ref]{sec} मध्ये आपण यावर
भर दिला होता की $\Real$ म्हणजे डेडेकिंड छेदांचा संच ही व्याख्या सोयीची
संकेतजनक व्याख्या न मानता \emph{शोध} मानणे चुकीचे ठरेल. डेडेकिंड छेद
वास्तव संख्यांच्या प्रमुख गणितीय गुणधर्मांचे उदाहरण देतात, हे निर्णायक
निरीक्षण आहे. येथेही तसेच: \emph{कोणतीही} डेडेकिंड बीजसंरचना नैसर्गिक
संख्यांच्या प्रमुख गणितीय गुणधर्मांचे उदाहरण देते, हे निर्णायक निरीक्षण
आहे. (खरे तर सर्व डेडेकिंड बीजसंरचना परस्पर \emph{समरूपी} आहेत, हे
सिद्ध करून डेडेकिंड यांनी हा मुद्दा ठामपणे दाखवला
(\citeyear[Theorems
132--3]{Dedekind1888}). त्यामुळे आजचे अनेक
``संरचनावादी'' डेडेकिंड यांना पूर्वसूरी मानतात, यात आश्चर्य नाही.)

 शिवाय आपण नैसर्गिक संख्यांचा सिद्धान्त अशा अनौपचारिक साध्या
 संचसिद्धान्तात कसा अंतःस्थापित करायचा हे दाखवले आहे, जो स्वतः अजूनही
 बराच अनौपचारिक आहे, पण नैसर्गिक संख्या आयत्या गृहीत धरत नाही असे
 (वरवर तरी) दिसते. त्यामुळे डेडेकिंड यांचा स्वतःचा महत्त्वाकांक्षी
 प्रकल्प प्रत्यक्षात आणण्याच्या मार्गावर आपण असू शकतो. त्यांनी तो असा
 स्पष्ट केला:
\begin{quote}
	विज्ञानात सिद्ध करण्याजोगी कोणतीही गोष्ट सिद्धतेशिवाय मान्य करू नये.
	ही मागणी रास्त वाटत असली, तरी अगदी साध्या विज्ञानाचा पाया घालण्याच्या
	सर्वांत अलीकडील पद्धतींमध्येही ती पूर्ण झाली आहे असे मला मानता येत
	नाही; म्हणजे, संख्यासिद्धान्ताशी संबंधित तर्कशास्त्राच्या त्या भागातही.
	अंकगणित (बीजगणित, विश्लेषण) हे तर्कशास्त्राचाच केवळ एक भाग आहे असे
	म्हणताना माझा आशय असा की संख्या ही संकल्पना अवकाश आणि काळ यांच्या
	कल्पना किंवा अंतःप्रज्ञा यांपासून पूर्णपणे स्वतंत्र आहे---ती विचारांच्या
	शुद्ध नियमांचे थेट फलित आहे असेच मी मानतो.
	\citep[preface]{Dedekind1888}
\end{quote}
डेडेकिंड यांची धाडसी कल्पना अशी आहे. आपण नुकतेच पाहिले की केवळ
(अनौपचारिक) संचसिद्धान्त वापरून नैसर्गिक संख्या कशा घडवायच्या.
\olref[sfr][arith][]{chap} मध्ये आपण पाहिले की नैसर्गिक संख्या आणि काही
संचसिद्धान्त दिला असता वास्तव संख्यांची रचना कशी करायची. त्यामुळे कदाचित
``अंकगणित (बीजगणित, विश्लेषण)'' हे ``तर्कशास्त्राचाच केवळ एक भाग'' ठरते
(डेडेकिंड यांच्या ``तर्कशास्त्र'' या शब्दाच्या विस्तारित अर्थाने).

ही झाली कल्पना. पण क्षणभर थांबा. डेडेकिंड बीजसंरचनेची आपली रचना
(म्हणजे नैसर्गिक संख्यांचे आपले पर्यायी प्रतिरूप) डेडेकिंड-अनंत संचाच्या
अस्तित्वावर अवलंबून आहे. (मागे \olref[sfr][infinite][dedekind]{thm:DedekindInfiniteAlgebra}
पाहा.) डेडेकिंड-अनंत संचाचे अस्तित्व ``तर्कशास्त्र'' किंवा ``विचारांचे
शुद्ध नियम'' यांच्या साहाय्याने प्रस्थापित करता आले नाही, तर हा प्रकल्प
अडून पडतो.

मग डेडेकिंड-अनंत संचाचे अस्तित्व ``विचारांच्या शुद्ध नियमांनी''
प्रस्थापित \emph{करता येते का}? डेडेकिंड यांचा प्रयत्न असा होता:
\begin{quote}
  माझे स्वतःचे विचारविश्व, म्हणजे माझ्या विचारांचे विषय होऊ शकणाऱ्या सर्व
  वस्तूंची समग्रता $S$, अनंत आहे. कारण $s$ हा~$S$ चा घटक दर्शवत असेल,
  तर $s$ हा माझ्या विचाराचा विषय होऊ शकतो हा विचार $s'$ स्वतः~$S$ चा
  घटक आहे. याला घटक~$s$ ची प्रतिमा $\phi(s)$ मानले, तर
  \dots~$S$ हा [डेडेकिंडच्या अर्थाने] अनंत आहे, हेच सिद्ध करायचे होते.
	\citep[\S66]{Dedekind1888}
\end{quote}
मुख्यतः अत्यंत काटेकोर गणितीय सिद्धतांनी भरलेल्या ग्रंथाच्या मध्यभागी
अशी गोष्ट सापडणे खरोखरच चकित करणारे आहे. येथे दोन टिप्पण्या कराव्याशा
वाटतात.

पहिली: या ``सिद्धतेला'' आज आपण ``गणितीय'' स्वरूप म्हणून ओळखू असे स्वरूप
जवळजवळ नाहीच. ती मानसशास्त्रीय वस्तूंबद्दल (विचारांबद्दल) बोलते आणि
त्यातही केवळ \emph{संभाव्य} विचारांबद्दल बोलते.

दुसरी: निदान आपण डेडेकिंड बीजसंरचना ज्या रीतीने मांडल्या आहेत त्या
रीतीनुसार, या ``सिद्धतेत'' एक सरळ तांत्रिक उणीव आहे. डेडेकिंड यांचा
युक्तिवाद यशस्वी झाला, तरी तो केवळ अनंत वस्तू (विशेषतः अनंत विचार)
अस्तित्वात आहेत एवढेच प्रस्थापित करतो. पण~$S$ ला अनेक !!{element}s असलेला
एकच \emph{संच} का मानावे,~$S$ ला केवळ \emph{काही वस्तू} (अनेकवचनात) का मानू
नये, याचे कारणही डेडेकिंड यांनी द्यायला हवे.

डेडेकिंड यांना येथे उणीव दिसली नाही, यावरून कदाचित त्यांनी वापरलेला
``समग्रता'' हा शब्द \emph{आपण} वापरत असलेल्या ``संच'' या शब्दाशी
तंतोतंत जुळत नाही असे सूचित होते.\footnote{तसे तो जुळत नव्हता असे
मानण्याची इतर कारणेही आपल्याकडे आहेत; पाहा \citet[p.~23]{Potter2004}.}
पण हे फार आश्चर्यकारक ठरणार नाही. मागील दोन प्रकरणांत आपण राबवलेला
प्रकल्प---नैसर्गिक संख्यांची ``रचना'' आणि त्यांच्यापासून पूर्णांक,
वास्तव संख्या व परिमेय संख्या यांची ``रचना''---संपूर्णपणे अनौपचारिक
रीतीने राबवला आहे. नेमके कोणते संच कधी अस्तित्वात असतात याविषयी अनेक
सामान्य तत्त्वे न मांडता, गरज पडेल तेव्हा आपण हा संच किंवा तो संच आयता
मानला. पण काही सामान्य तत्त्वांची आपल्याला \emph{गरज} आहे हे माहीत आहे;
नाहीतर आपल्याला रसेलच्या विरोधापत्तीला सामोरे जावे लागेल.

 आता अनौपचारिकतेच्या या मर्यादा ओलांडण्याची वेळ आली आहे.

\end{document}
